\subsection{克利福德代数的定义与泛性质.}

定义 1. --- 我们称模 E 的张量代数 T(E) 对由形如 \(x \otimes x - Q(x)\) . 1（\(x \in E\)）的元素生成的双边理想（记作 I(Q)）的商代数为 Q 的克利福德代数，并用 C(Q) 记之。

  我们将用 \(\rho_{\mathbb{Q}}\)（当不致混淆时也简记作 \(\rho\)）记由 E 到 T(E) 中的典范映射与 T(E) 到 C(Q) 上的典范映射 \(\sigma\) 复合而得的从 E 到 C(Q) 中的映射；映射 \(\rho_{\mathbb{Q}}\) 称为典范的。注意 C(Q) 由 \(\rho_{\mathbb{Q}}(E)\) 生成，且对 \(x \in E\)，有

\[
\rho (x) ^ {2} = \mathrm{Q} (x). 1;\tag{1}
\]

从而，把 \(x\) 换成 \(x + y\)（ \(x, y\) 属于 E）

\[
\rho (x) \rho (y) + \rho (y) \rho (x) = \Phi (x, y). 1\tag{2}
\]

例. --- 若 E 承认一个由单个元素 \(e\) 组成的基，则 T(E) 同构于多项式代数 A{[}X{]}，而 C(Q) 是 A 的一个二次扩张，以 \((1, u)\) 为基，其中 \(u\) 是元素 \(u = \rho(e)\)，它满足 \(u^2 = Q(e)\)。

  用 \(T^{h}\) 记 h 次张量幂 \(\otimes E\)（在 \(\mathrm{T}(\mathrm{E})\) 中），并设 \(T^{+}\)（相应地，\(T^{-}\)）是 h 为偶数（相应地，奇数）的那些 \(T^{h}\) 之和。

  由于 T(E) 是 \(T^{+}\) 与 \(T^{-}\) 的直和，且 I(Q) 由 \(T^{+}\) 的元素生成，故 I(Q) 是 \(T^{+} \cap I(Q)\) 与 \(T^{-} \cap I(Q)\) 的直和，而 C(Q) 是两个子模 \(C^{+}(Q) = \sigma(T^{+})\) 与 \(C^{-}(Q) = \sigma(T^{-})\)（也记作 \(C^{+}\) 与 \(C^{-}\)）的直和。\(C^{+}\) 的元素将称为偶的（相应地，奇的）。我们有关系式

\[
\mathrm{C} ^ {+} \mathrm{C} ^ {+} \subset \mathrm{C} ^ {+}, \quad \mathrm{C} ^ {+} \mathrm{C} ^ {-} \subset \mathrm{C} ^ {-}, \quad \mathrm{C} ^ {-} \mathrm{C} ^ {+} \subset \mathrm{C} ^ {-}, \quad \mathrm{C} ^ {-} \mathrm{C} ^ {-} \subset \mathrm{C} ^ {+}.\tag{3}
\]

  特别地，\(C^{+}\) 是 \(C(Q)\) 的一个子代数。

命题 1. --- 设 f 是 E 到 A 上一个代数 D 中的线性映射，使得 \(f(x)^{2} = \mathrm{Q}(x)\) . 1 对每个 \(x \in \mathbf{E}\) 成立。存在一个且只有一个 C(Q) 到 D 中的同态 \(\bar{f}\)，使得 \(f = \bar{f} \circ \rho_{\mathbb{Q}}\)。

  \(\bar{f}\) 的唯一性由 C(Q) 由 \(\rho_{\mathbb{Q}}(\mathrm{E})\) 生成这一事实得出。设 \(h\) 是扩张 \(f\) 的 T(E) 到 D 中的唯一同态（\(h\) 由 \(h(x_1 \otimes \cdots \otimes x_n) = f(x_1) \ldots f(x_n)\) 定义）。我们有

\[
h (x \otimes x - \mathrm{Q} (x). 1) = (f (x) ^ {2} - \mathrm{Q} (x)). 1 = 0,
\]

  因而 h 在 I(Q) 上为零，并通过过渡到商定义所求的同态 \(\bar{f}\)。

  命题 1 表明 C(Q) 是一个泛映射问题的解（《集合论》第四章 § 3 第 1 小节）。

  特别地，取 D 为 C(Q) 的反代数，取 \(f\) 为映射 \(\rho\)；命题 1 蕴涵存在一个且只有一个 C(Q) 的反自同构 \(\beta\)，它在 \(\rho(E)\) 上的限制是恒等映射；它称为 C(Q) 的主反自同构。显然 \(\beta^{2} = 1\)。

  另一方面，设 Q' 是 A-模 E' 上的一个二次型，f 是 E 到 E' 中的线性映射，使得 Q'∘f = Q。我们有 \(\rho_{\mathbb{Q}'}(f(x))^{2} = \mathrm{Q}'(f(x)).1 = \mathrm{Q}(x).1\)，因而存在一个且只有一个从 C(Q) 到 C(Q') 中的同态 C(f)，使得 C(f)∘ρq = ρq'∘f。若 f 是恒等映射，则 C(f) 是恒等映射；若 Q'' 是 A-模 E'' 上的一个二次型，g 是 E' 到 E'' 中的线性映射，使得 Q''∘g = Q'，则我们有 C(g∘f) = C(g)∘C(f)。当 E' 是 E 的子模且 f 是 E' 到 E 中的典范单射（于是 Q' 是 Q 在 E' 上的限制）时，我们称 C(f) 为 C(Q') 到 C(Q) 中的典范同态。

  特别地，取 Q' = Q 并取 f 为映射 \(x \to -x\)；我们看出存在一个且只有一个 C(Q) 的自同构 \(\alpha\)，使得 \(\alpha \circ \rho = -\rho\)；它称为 C(Q) 的主自同构。显然 \(\alpha^{2} = 1\)，且 \(\alpha\) 在 C⁺（相应地，C⁻）上的限制是恒等映射（相应地，映射 \(u \to -u\)）。

命题 2. --- 设 A' 是一个交换环，φ 是 A 到 A' 中的同态，Q' 是由 Q 经纯量扩张得到的 E' = A' ⊗\_A E 上的二次型（§ 3 第 4 小节命题 3）。存在一个且只有一个代数 A' ⊗\_A C(Q) 到 C(Q') 上的同构 j，使得对一切 x ∈ E 有 j (1 ⊗ ρ\_Q(x)) = ρ\_Q′(1 ⊗ x)。

  只需证明代数 \(C' = A' \otimes C(Q)\) 与映射 \(1 \otimes_{\rho_{Q}}\)（从 \(E'\) 到 \(C'\) 中）构成与 \(C(Q')\) 和 \(\rho_{Q'}\) 相同的泛问题的一个解。事实上，设 \(D'\) 是 \(A'\) 上一个代数，\(f'\) 是 \(E'\) 到 \(D'\) 中的一个 A'-线性映射，使得 \(f'(x')^{2} = Q'(x')\) . 1 对每个 \(x' \in E'\) 成立。映射 \(g : x \to f'(1 \otimes x)\)（从 E 到 \(D'\) 中，后者经由同态 \(\varphi\) 看作 A-模）是 A-线性的，且 \(g(x)^{2} = Q'(1 \otimes x)\) . 1 = Q(x). 1 对一切 \(x \in E\) 成立。因此存在一个且只有一个 A-同态 \(\bar{g}\)（从 C(Q) 到 \(D'\) 中），使得 \(\bar{g}(\rho_{Q}(x)) = f'(1 \otimes x)\)。从而存在一个且只有一个 A'-同态 \(\overline{f'}\)（从 C' 到 \(D'\) 中），使得 \(\overline{f}'(1 \otimes \rho_{Q}(x)) = f'(1 \otimes x)\) 对一切 \(x \in E\) 成立；由线性性得出，\(\overline{f}'((1 \otimes \rho_{Q})(x')) = f'(x')\) 对一切 \(x' \in E'\) 成立。证毕。

